% Problem record op_378472b2da3c533d. This TeX file is derived from
% database/problems_json/op_378472b2da3c533d.json by scripts/sync-tex.mjs; edit the JSON record
% and rerun the script rather than editing this file.
\section{Average-case approximation hardness of random Ising partition functions}
\label{sec:op_378472b2da3c533d}

\paragraph{Problem.}
Is it $\#\mathrm{P}$-hard to approximate $|Z_R|^2$ to relative multiplicative
error $a+o(1)$ on a $b$ fraction of random Ising instances?
Here $a>0$ and $0<b\leq1$ are constant parameters independent of the number
of vertices $n$, specifying the relative error and instance fraction.

Take the complete graph on $n$ vertices. Choose each edge weight $w_{ij}$
and each vertex weight $v_k$ independently and uniformly from
$\{0,\ldots,7\}$. With $\omega=e^{i\pi/8}$, define
\begin{equation}
  Z_R=\sum_{z\in\{-1,1\}^n}
    \omega^{\sum_{i<j}w_{ij}z_i z_j+\sum_{k=1}^n v_k z_k}.
  \label{eq:bms-ising-partition}
\end{equation}
The target is the squared modulus of Eq.~\eqref{eq:bms-ising-partition}.
An estimate $\widetilde Q_R$ is required to satisfy
\begin{equation}
  \bigl|\widetilde Q_R-|Z_R|^2\bigr|
    \leq (a+o(1))|Z_R|^2.
  \label{eq:bms-ising-approximation}
\end{equation}
The $b$ fraction in the question is measured over the random vertex and edge
weights for which Eq.~\eqref{eq:bms-ising-approximation} holds; $o(1)$ tends
to zero as $n\to\infty$.

\subsection*{Status}

Unsolved

\subsection*{Source}

Conjecture 2 of Bremner, Montanaro, and Shepherd, on page 2 of arXiv v2,
states this average-case hardness conjecture with $a=1/4$ and $b=1/24$
\sourcecite{ref:bms-ising-source}{BMS16}.
The present formulation replaces those two numerical constants by the
parameters $a$ and $b$, as requested by the contributor, and retains the
original $o(1)$ term, random-instance distribution, and squared-modulus target.
The parameterized formulation is not a verbatim claim of the paper.

\subsection*{Progress}

\begin{itemize}
  \item None
\end{itemize}

\subsection*{References}

\begin{enumerate}[label={},leftmargin=4.8em,labelsep=0.5em]
  \footnotesize
  \item[\textup{[BMS16]}]\label{ref:bms-ising-source}
  M. J. Bremner, A. Montanaro, and D. J. Shepherd,
"Average-case complexity versus approximate simulation of commuting quantum computations,"
\emph{Physical Review Letters} \textbf{117}, 080501 (2016).
\href{https://doi.org/10.1103/PhysRevLett.117.080501}{doi:10.1103/PhysRevLett.117.080501};
\href{https://arxiv.org/abs/1504.07999v2}{arXiv:1504.07999v2}.
\end{enumerate}

\subsection*{Comment}

The remaining task is an average-case hardness result for the specified
random-weight distribution and relative-error guarantee. The pair $(a,b)$
indexes a family of questions; hardness for every pair is not asserted.
The source's original parameter choice is $(a,b)=(1/4,1/24)$.
The instance fraction $b$ concerns inputs, not an algorithm's internal
success probability.

\subsection*{Field}

Quantum algorithm

\subsection*{Topic}

Computational complexity; Quantum supremacy; IQP sampling

\subsection*{ID}

\texttt{op\_378472b2da3c533d}
